%! Principal Component Analysis — Unit 7, Lesson 7.3 — Guided Notes (Answer Key)
\documentclass[10pt]{article}
\usepackage{linalg-article}
\usepackage{linalg-key}   % pulls in -boxes; do NOT also load -boxes
\usetikzlibrary{arrows.meta}
\newcommand{\TallMath}[1]{$\displaystyle #1\rule[-1.4em]{0pt}{3.2em}$}
\newcommand{\uu}{\mathbf{u}}
\newcommand{\vv}{\mathbf{v}}
\newcommand{\ee}{\mathbf{e}}
\newcommand{\zero}{\mathbf{0}}
\newcommand{\T}{^{\top}}

\begin{document}

\pageheader{Unit 7, Lesson 7.3}{Guided Notes --- Answer Key}
\namedateperiod

\vspace{0.12in}

\begin{objectivebox}
\textbf{By the end of this lesson, I will be able to\dots}
\begin{itemize}\setlength{\itemsep}{2pt}
  \item center a data table and form $A\T A$, then read its eigenvectors as the principal components $\vv_1\perp\vv_2$;
  \item measure the fraction of the spread the first component captures, and interpret what PC1 means for the data;
  \item explain why keeping the top few components summarizes data in far fewer numbers.
\end{itemize}
\end{objectivebox}

\vspace{0.06in}

\begin{vocabbox}
\small
Fill in each term as we build it together.
\vspace{2pt}
\vocabans{Centered data (subtract each column's mean --- the cloud sits at the origin)}{a data table with each column's mean subtracted, so the data cloud is centered at the origin}
\vocabans{Principal component (an eigenvector of $A\T A$ --- a direction the data spreads along)}{an eigenvector of $A\T A$; the perpendicular directions the data spreads along, ordered biggest $\sigma$ first}
\vocabans{PC1 (the direction of \emph{greatest} spread --- the best single direction through the cloud)}{the eigenvector with the largest eigenvalue --- the direction of most spread, the best-fit line through the cloud}
\end{vocabbox}

\begin{hookbox}
Four students each took two quizzes; the scatter of their $(\text{quiz 1},\text{quiz 2})$ points leans along a
diagonal --- do well on one, do well on the other.
\begin{itemize}\setlength{\itemsep}{1pt}
  \item If you had to describe each student with a \emph{single} number instead of two, along what direction would you measure? \ansline{Along the diagonal the cloud leans on --- the ``both scores rise together'' direction (this is PC1).}
  \item Roughly how much of the story would that one direction capture --- and what would you lose? \ansline{Most of it; you lose only the small spread \emph{across} the diagonal (how much better one quiz was than the other).}
\end{itemize}
\end{hookbox}

\begin{notesbox}{1. Data is a Matrix --- First, Center It}
Write the data as a matrix $A$: \textbf{one row per data point}, \textbf{one column per feature}. Our four
students, scores $(\text{quiz 1},\text{quiz 2})$, are $(5,6),(6,5),(3,2),(2,3)$.

Before we look for directions, we \textbf{center} the data: subtract each column's mean so the cloud sits at
the origin. Then a direction describes how the data \emph{spreads}, not where it happens to sit.
\[
  \text{mean of each column}=(\,{\color{keyred}\mathbf{4}},{\color{keyred}\mathbf{4}}\,)\ \Rightarrow\
  \text{centered rows: } {\color{keyred}\mathbf{(1,2),\ (2,1),\ (-1,-2),\ (-2,-1)}}.
\]
From now on $A$ is the \emph{centered} matrix
$A=\begin{bsmallmatrix} 1 & 2\\ 2 & 1\\ -1 & -2\\ -2 & -1 \end{bsmallmatrix}$.
\end{notesbox}

\begin{notesbox}{2. The Principal Components are the Eigenvectors of $A\T A$}
This is exactly the Lesson~7.1 recipe. Form the symmetric matrix $A\T A$ (each entry is a dot product of two
columns of $A$ --- the spreads and overlaps of the features):
\[
  A\T A=\begin{bmatrix}{\color{keyred}\mathbf{10}} & {\color{keyred}\mathbf{8}}\\[2pt]{\color{keyred}\mathbf{8}} & {\color{keyred}\mathbf{10}}\end{bmatrix}
  =\begin{bmatrix} 10 & 8\\ 8 & 10 \end{bmatrix}\quad(\text{from the Warm-Up}).
\]
Because $A\T A$ is \emph{symmetric}, its eigenvectors are \textbf{perpendicular} (Lesson~7.1) --- they are the
natural axes of the data cloud, called the \textbf{principal components}. From the Warm-Up, $\lambda=18,2$, so
$\sigma_1=\sqrt{18},\ \sigma_2=\sqrt2$. Solve $(A\T A-\lambda I)\vv=\zero$ for each:
\[
  \lambda=18:\ \vv_1=\tfrac1{\sqrt2}\begin{bmatrix}{\color{keyred}\mathbf{1}}\\[2pt]{\color{keyred}\mathbf{1}}\end{bmatrix}
  \ (\textbf{PC1}),\qquad
  \lambda=2:\ \vv_2=\tfrac1{\sqrt2}\begin{bmatrix}{\color{keyred}\mathbf{1}}\\[2pt]{\color{keyred}\mathbf{-1}}\end{bmatrix}
  \ (\textbf{PC2}).
\]
The bigger eigenvalue $\lambda_1=18$ belongs to \textbf{PC1} --- the direction the data spreads the most.
\end{notesbox}

\begin{notesbox}{3. Reading the Components: How Much Spread?}
\begin{minipage}[t]{0.56\linewidth}\vspace{0pt}
Each eigenvalue \emph{is} the spread along its direction. If you drop every centered point straight onto the
line through $\vv$ and add up the squared distances, you get exactly $\lambda=\sigma^2$. So
$\sigma_1^2=18$ is the spread along PC1 and $\sigma_2^2=2$ is the spread along PC2.

\vspace{3pt}
The \textbf{fraction of the total spread} carried by PC1 is
\[
  \frac{\sigma_1^2}{\sigma_1^2+\sigma_2^2}=\frac{18}{18+2}=\frac{18}{20}={\color{keyred}\mathbf{90}}\%.
\]
PC1 alone captures most of the cloud; PC2 holds the little that is left. (Statisticians divide these by
$n-1$ to call them \emph{variances} --- but the \emph{fraction} is the same either way.)
\end{minipage}\hfill
\begin{minipage}[t]{0.40\linewidth}\vspace{0pt}
\centering
\begin{tikzpicture}[scale=0.92, font=\footnotesize]
  % axes
  \draw[linegray, -{Latex[length=1.6mm]}] (-2.4,0)--(2.4,0);
  \draw[linegray, -{Latex[length=1.6mm]}] (0,-2.4)--(0,2.4);
  % centered data points (scaled 0.8)
  \foreach \x/\y in {0.8/1.6, 1.6/0.8, -0.8/-1.6, -1.6/-0.8}
    \fill[charcoal] (\x,\y) circle (2.1pt);
  % PC1 along (1,1): long burgundy
  \draw[burgundy, -{Latex[length=2.4mm]}, very thick] (0,0)--(1.9,1.9);
  \node[burgundy] at (1.5,2.15) {\textbf{PC1} ($\sigma_1$)};
  % PC2 along (1,-1): short royalblue
  \draw[royalblue, -{Latex[length=2.4mm]}, very thick] (0,0)--(0.8,-0.8);
  \node[royalblue] at (1.55,-0.95) {\textbf{PC2} ($\sigma_2$)};
\end{tikzpicture}\\[-1pt]
{\scriptsize The centered cloud spreads far along PC1, little along the perpendicular PC2.}
\end{minipage}
\end{notesbox}

\begin{notesbox}{4. What the Components \emph{Mean} --- and Dimension Reduction}
\textbf{Interpret the directions.} $\vv_1=\tfrac1{\sqrt2}\begin{bsmallmatrix} 1\\1 \end{bsmallmatrix}$ points
where \emph{both} quizzes rise together --- an ``overall score'' axis. $\vv_2=\tfrac1{\sqrt2}\begin{bsmallmatrix} 1\\-1 \end{bsmallmatrix}$
is the ``which quiz did you do better on'' axis. PC1 carries $90\%$ of the spread, so \emph{overall score} is
what really separates these students.

\vspace{2pt}
\textbf{Dimension reduction.} Summarize each student by their \emph{coordinate along PC1} alone --- one number
instead of two --- and you keep $90\%$ of the information. That is PCA's payoff: with $100$ features, keeping
the top $2$ or $3$ components replaces $100$ numbers per point with a handful, losing very little.

\vspace{2pt}
\textbf{This is Lesson 7.2, on data.} Compressing an image kept the biggest-$\sigma$ \emph{layers}; PCA keeps
the biggest-$\sigma$ \emph{directions}. Same SVD, same ``keep the biggest'' rule --- pixels became data.

\vspace{2pt}
\textbf{The road ahead.} \textbf{7.4} the victory of orthogonality --- why perpendicular directions are what
make the SVD, compression, and PCA all work.
\end{notesbox}

\begin{practicebox}
Show your work.
\begin{enumerate}[leftmargin=1.6em, itemsep=4pt]
  \item A data set has centered rows $(2,2),(1,1),(-1,-1),(-2,-2)$ --- all on the line $y=x$. Form $A\T A$ and find its eigenvalues. What is $\sigma_2$, and what does that say about the cloud? \ansline{$A\T A=\begin{bsmallmatrix} 10 & 10\\ 10 & 10 \end{bsmallmatrix}$; $\lambda=20,0$, so $\sigma_2=0$. The cloud has \emph{no} spread off PC1 --- it lies exactly on one line, so a single principal component describes it perfectly.}
  \item For the notes data, PC1 $=\tfrac1{\sqrt2}(1,1)$ carries $90\%$ of the spread. In one sentence, what does ``moving along PC1'' mean for a student? \ansline{Both quiz scores go up (or down) together --- PC1 measures overall performance, the main way these students differ.}
  \item Why must we \emph{center} the data before finding principal components? \ansline{Centering puts the cloud at the origin so the eigenvectors of $A\T A$ describe how the data \emph{spreads}; on un-centered data the top direction just points toward the mean, not the spread.}
\end{enumerate}
\end{practicebox}

\begin{teachernote}
The new setup move is \emph{centering} (Section~1): subtract each column's mean so directions measure spread,
not location. Section~2 is the punchline that ties PCA to the whole unit --- the principal components are just
the eigenvectors of the symmetric $A\T A$ (the Lesson~7.1 recipe), so they are automatically perpendicular
and ordered by $\sigma$. The spine (four students, two quiz scores) is built so $A\T A=\begin{bsmallmatrix} 10 & 8\\ 8 & 10 \end{bsmallmatrix}$
gives clean $\lambda=18,2$ and the $\pm45^\circ$ components $\tfrac1{\sqrt2}(1,\pm1)$. Section~3's key idea:
each eigenvalue \emph{is} the spread along its direction (sum of squared projections $=\vv\T A\T A\vv=\lambda$),
so PC1 holds $18/20=90\%$. Section~4 lands the meaning --- PC1 $=$ ``overall score,'' PC2 $=$ ``which quiz was
better'' --- and dimension reduction: keep the top components, drop the rest, exactly like keeping the biggest
image layers in 7.2. Common slips: skipping the centering step; forgetting to \emph{square} $\sigma$ for the
fraction; calling a component a data point instead of a direction.
\end{teachernote}

\end{document}
